
\setlist[itemize,1]{label=\textbullet}

\section{Démonstrations des propriétés du parser}

On souhaite démontrer les propriétés de correction et de terminaison du parser on introduit pour cela les notations suivantes:

Soit $\mathcal{P} = (\mathcal{O}, \mathcal{A}, Pred, S_0, \mathcal{G})$ un problème PDDL avec:
\begin{itemize}
    \item $\mathcal{O}$: l'ensemble des objets du problème
    \item $\mathcal{A}$: l'ensemble des actions du problème avec pour tout $a\in \mathcal{A}$ la définition du tuple:  $$(param(a),pre(a),eff^+(a),eff^-(a))$$
    \item $Pred$: l'ensemble des prédicats déclarés dans le domaine de $\mathcal{P}$
    \item $S_0$: l'ensemble des prédicats vérifiés dans l'état initial (hypothèse de monde clos), $S_0 \in 2^{\mathds{F}(\mathcal{P})}$
    \item $\mathcal{G}$: formule de l'objectif construite sur $\mathds{F}(\mathcal{P})$, décrit l'ensemble des conditions qui doivent être respecté pour considérer un état comme final.
    \item $\mathds{F}(\mathcal{P}) = \{ p(o_1,\dots,o_k) \mid p\in Pred, (o_1, \dots,o_k) \in \mathcal{O}^k \}$ (ensemble fini car $Pred$ et $\mathcal{O}$ sont finis)
    \item $\mathds{E}(\mathcal{P}) = \{ a(o_1,\dots, o_k) \mid a \in \mathcal{A}, (o_1, \dots,o_k) \in \mathcal{O}^k, k = \lvert param(a) \rvert \}$
\end{itemize}

On introduit $\langle \mathds{S}^*, \tau^*\rangle$ le LTS cible associé à $\mathcal{P}$ avec: 
$$\mathds{S} = \{S  \mid S \in 2^{\mathds{F}(\mathcal{P})} \text{ accessible depuis }S_0\}$$
\begin{equation*}
    \begin{aligned}
    \tau^* &= \{(S, \{e\}, S') \mid S \in \mathbb{S}^*, e \in \mathbb{E}(\mathcal{P}), pre(e) \text{ satisfaite dans } S \text{ et } S' = (S \setminus eff^-(e)) \cup eff^+(e)\} \\
    &\cup \{(S, \emptyset, S) \mid S \in \mathbb{S}^+, S \models \mathcal{G}\}
    \end{aligned}
\end{equation*}

et $T(\mathcal{P}) = \langle S, \tau \rangle $ avec $T$ notre processus (parser) et donc $T(\mathcal{P})$ le LTS obtenu en l'appliquant au problème PDDL $\mathcal{P}$.


\subsection{Correction}
Pour démontrer la correction on souhaite montrer dans un premier temps que le LTS obtenu par notre parser est correct. C'est à dire que $\mathds{S} \subseteq 2 ^{\mathds{F}(\mathcal{P})}$ et $\tau \subseteq \mathds{S} \times 2 ^{\mathds{E}(\mathcal{P})} \times \mathds{S} $. Par construction $S_0 \in \mathds{S}$ et $S_0\in 2 ^{\mathds{F}(\mathcal{P})}$. De plus pour tout $S\in \mathds{S}$ et pour tout $e \in \mathds{E}(\mathcal{P})$ si $pre(e)$ est vérifié dans $S$ on construit l'état résultant de l'application de $e$ dans $S$ de la façon suivante $S' = (S \setminus eff^-(e)) \cup eff^+(e)$ et $eff^+(e),eff^-(e) \in 2 ^{\mathds{F}(\mathcal{P})}$ donc $S'\in 2 ^{\mathds{F}(\mathcal{P})}$. Et on ajoute $(S,\{e\}, S')$ à $\tau$. Cette procédure est la seule possible pour ajouter des états à $\mathds{S}$ donc $\mathds{S} \subseteq 2 ^{\mathds{F}(\mathcal{P})}$. 
Pour $\tau$ on ajoute uniquement en plus les arcs réflexifs qui appartiennent bien à $\mathds{S} \times 2 ^{\mathds{E}(\mathcal{P})} \times \mathds{S}$ donc on a bien $\tau \subseteq \mathds{S} \times 2 ^{\mathds{E}(\mathcal{P})} \times \mathds{S} $. On en conclut que $T(\mathcal{P})$ est un LTS correct.

Dans un second temps on souhaite montrer $T(\mathcal{P}) = \langle \mathds{S}, \tau \rangle$ avec $S = \mathds{S}^*$ et $\tau = \tau^*$.
\begin{itemize}
    \item $\mathds{S} \subseteq \mathds{S}^*$: pour tout $S\in \mathds{S}$, $S$ a été atteint par BFS depuis $S_0$ en aplique des événements dont les préconditions sont satisfaites donc $S$ est accessible depuis $S_0$ i.e $\mathds{S} \subseteq \mathds{S}^*$
    \item $\tau \subseteq \tau^*$: pour toutes transitions $(S, \{e\}, S')\in \tau$,  $pred(e)$ est vérifié dans $S$ et $S' = (S \setminus eff^-(e)) \cup eff^+(e)$ donc $(S, \{e\}, S') \in \tau^*$ et les arêtes de la forme $(S, \emptyset, S)\in \tau$ sont ajoutées si et seulement si  $S \models \mathcal{G}$ donc $(S, \emptyset, S)\in \tau^*$.
    \item $\mathds{S} \supseteq \mathds{S}^* $ et $\tau \supseteq \tau^*$ par induction: 
    \begin{itemize}
        \item Base: $S_0 \in \mathds{S}$ car l'algorithme BFS est initialisé avec $S_0$
        \item Induction: Soit $S\in \mathds{S}^*$ accessible depuis $S_0$ par définition. Soit $e \in \mathds{E}(\mathcal{P})$ tel que $pre(e)$ est vérifié dans $S$ et soit $S' = (S \setminus eff^-(e)) \cup eff^+(e)$ par définition $S'\in \mathds{S}^*$ et $(S,\{e\}, S') \in \tau*$. Lorsque le processus $T$ va traiter $S$ il va tester l'ensemble des événements et ajouter $(S,\{e\}, S')$ à $\tau$ et $S'$ à $\mathds{S}$ ainsi qu'à la liste d'états à traiter. On a bien $\mathds{S} \supseteq \mathds{S}^* $ et $\tau \supseteq \tau^*$ .
        \end{itemize}
 \end{itemize}

\subsection{Terminaison}
Le processus $T$ est basé sur un algorithme BFS donc l'espace d'exploration est fini et borné par le nombre d'états possible égal à $2^{\mid \mathds{F}(\mathcal{P}) \mid}$. Les états ne sont traité qu'une fois avec l'usage de la liste "visited".